% Results-pending abstract: investigation language is intentional.
% Add measured findings only after completing the corresponding experiments.
\begin{abstract}
Multi-teacher on-policy distillation (MOPD) aims to consolidate specialized
capabilities in a single student. We examine three aspects of its
optimization: token balancing, update sparsity, and supervision
density. Starting from the MOPD learning objective, we first show how
loss averaging determines both domain weights and the relative weight
of long and short responses. A finite-batch gradient decomposition
separates these effects and motivates matched experiments on capability
balance. Second, we investigate whether parameter changes concentrate
on small, overlapping subsets when different teachers act on the same
MOPD student, and whether more different teacher distributions correspond
to more different subsets. We measure optimizer steps and cumulative
parameter changes separately, using raw gradients to help explain the
observed updates. Third, we compare sampled-token policy gradients with
teacher top-$k$ supervision in single-teacher and joint training,
using full-vocabulary reverse-KL updates at shared prefixes as a
local reference. The planned experiments pair update concentration
with per-domain capability and account for sampling variability,
training exposure, and compute. Each question is developed together
with its experimental design, without presuming that sparse updates,
low teacher overlap, or sampled supervision improve capability
integration. Empirical outcomes remain to be measured.
\end{abstract}
